% =============================================================================
% Resumen en castellano.
% =============================================================================

\thispagestyle{empty}

\noindent\textbf{\tfgtitulo}\\
\tfgsubtitulo.\\[0.4cm]

\noindent\tfgautor\\[0.6cm]

\noindent\textbf{Palabras clave}: detección de drones, radiofrecuencia, SDR, SigMF, espectrogramas, STFT, YOLOv8, detección de objetos, OcuSync, Bluetooth, WiFi, FHSS, ResNet18.

\vspace{0.6cm}
\section*{Resumen}

Este Trabajo Fin de Grado aborda la detección de drones comerciales mediante el análisis de señales de radiofrecuencia capturadas con receptores de radio definida por software. El trabajo se centra en la banda ISM de \SI{2.4}{\giga\hertz}, donde los enlaces de comunicación de drones comerciales coexisten con tecnologías de uso masivo como WiFi y Bluetooth. En este contexto, el problema no consiste únicamente en detectar energía en el espectro, sino en localizar y clasificar estructuras tiempo-frecuencia compatibles con la firma del dron frente a otras señales presentes en la misma banda.

La metodología propuesta se reformula como un sistema en dos etapas. En la primera etapa, las muestras IQ capturadas con USRP se almacenan en formato SigMF, se segmentan en ventanas de \SI{0.1}{\second} y se transforman en espectrogramas mediante la transformada de Fourier de tiempo corto. Sobre dichos espectrogramas se entrena un detector de objetos YOLOv8 para localizar ráfagas o \emph{bursts} FHSS y clasificarlas como \emph{drone} o \emph{interference}. Este enfoque sustituye al clasificador global de imagen completa, ya que fuerza al modelo a aprender la morfología local de las ráfagas y reduce el riesgo de aprender sesgos contextuales de la captura. La segunda etapa, basada en una red ResNet18, identifica el modelo concreto de dron a partir de las regiones ya filtradas por la primera.

El resultado más relevante de la primera etapa es que el detector no presenta confusión cruzada entre las clases activas \emph{drone} e \emph{interference}: cuando emite una predicción positiva sobre una región del espectrograma, la clase asignada es correcta, y los errores se limitan a falsos negativos hacia \emph{background}. Aunque el recall a nivel de ráfaga individual es moderado (precisión global de 0,760, recall de 0,391 y mAP@50 de 0,388 en test, con mAP@50 por clase de 0,575 para \emph{drone} frente a 0,201 para \emph{interference}), la redundancia temporal del protocolo OcuSync ---decenas de ráfagas por ventana de \SI{0.1}{\second}--- permite confirmar la presencia del dron a nivel de imagen con un recall del \SI{91}{\percent}, manteniendo precisión y especificidad de 1,000, sin un solo falso positivo sobre las 600 imágenes de prueba.

La segunda etapa identifica el DJI Mini~3 ---plataforma de referencia del trabajo--- con un F1 de 0,992, prácticamente sin errores. Un barrido de robustez frente a ruido demuestra además que la decisión de la red se sustenta en la morfología espectro-temporal de las ráfagas y no en su nivel medio de potencia: a igual potencia inyectada, un ruido blanco que cubre todo el plano tiempo-frecuencia degrada el clasificador mucho antes que una interferencia de salto de frecuencia, que solo ocupa una pequeña fracción del plano. Finalmente, el detector entrenado únicamente con el DJI Mini~3 localiza, sin reentrenamiento alguno, las ráfagas de otros drones no vistos en el entrenamiento (DJI F450 y Mavic~Pro, con protocolo OcuSync~v1), lo que confirma que la característica aprendida es la morfología FHSS genérica y no la firma de un modelo concreto. Estos resultados se obtienen sobre capturas en condiciones controladas; su validación en entornos de despliegue más diversos se plantea como línea futura.

\cleardoublepage
